\documentclass[a4paper,10pt]{article}
\usepackage[utf8]{inputenc}
\usepackage{graphicx}
\usepackage{subfigure}
\usepackage{fancyhdr}
\usepackage{listings}
\usepackage{ gensymb }


%opening

\pagestyle{fancy} %eigener Seitenstil
\fancyhf{} %alle Kopf- und Fußzeilenfelder bereinigen
\fancyhead[L]{SEP Abgabegespräch SS2015/16} %Kopfzeile links
\fancyhead[C]{} %zentrierte Kopfzeile
\fancyhead[R]{\includegraphics[width=1cm]{../../template/tugraz_bw}} %Kopfzeile rechts
\renewcommand{\headrulewidth}{0.4pt} %obere Trennlinie
\fancyfoot[C]{} %Seitennummer
\renewcommand{\footrulewidth}{0.4pt} %untere Trennlinie


\lstset{language=C++, basicstyle=\ttfamily\tiny}  

\begin{document}

\section*{Rollenspiel-Helden-Verwaltung}

Schreibe ein einfaches Programm um Rollenspiel-Helden für ein Computerspiel zu verwalten. Die Helden werden mit Namen erzeugt und besitzen ein Level sowie Erfahrungspunkte (XP). Die Erfahrungspunkte fürs nächste Level können wie folgt berechnet werden:

\begin{lstlisting}[basicstyle=\footnotesize]
int calcXPFormula(int current_level)
{
  return exp(current_level/2) + 10;
}
\end{lstlisting}

\begin{description}
 \item[create $<$heroName$>$] Erstellt einen neuen Helden mit dem Namen $heroName$. Das Level des Helden wird auf $1$ initialisiert, und die aktuellen Erfahrungspunkte auf $0$.
 \item[show] Listet alle Helden mit Namen und Erfahrungspunkten auf. Das Format der Ausgabe ist ''$<$heroName$>$ (Level $<$currentLevel$>$, $<$currentXP$>$ / $<$XPToNextLevel$>$ XP)''
 \item[gainXP $<$XPValue$>$] Alle Helden bekommen zusätzliche $XPValue$ Erfahrungspunkte. Falls ein Held das nächste Level erreicht ($currentXP$ $\geq$ $XPToNextLevel$), wird ''$<$heroName$>$ leveled up!'' ausgegeben. Es sind außerdem mehrere Level Ups auf einmal möglich.
\end{description}
Bei Fehlerfällen (Falsche Parameter, kein Held existiert...) soll nichts ausgegeben werden.
Bei der Eingabe von \texttt{quit} oder EOF soll das Programm beendet werden.

\subsection*{Beispiel}
\begin{verbatim}
> create Siegfried
> show
Siegfried (Level 1, 0 / 11 XP)
> gainXP 11
Siegfried leveled up!
Siegfried (Level 2, 0 / 12 XP)
> create Arthos
> gainXP 1000
Siegfried leveled up!
Siegfried (Level 13, 25 / 413 XP)
Arthos leveled up!
Arthos (Level 13, 14 / 413 XP)
> quit
\end{verbatim}

\subsection*{Aufruf}
\texttt{./rpg-hero-management}

\subsection*{Bewertung}
Dokumentation \& Programmierstil,  Dokumentation \& Programmierstil, und Struktur \& Korrektheit (je 1 Punkt)
\newpage


\end{document}
